% =============================================================================
%  About the Authors
% =============================================================================
\cleardoublepage
\chapter*{About the Authors}
\addcontentsline{toc}{chapter}{About the Authors}
\thispagestyle{plain}
\markboth{About the Authors}{}

\section*{\textcolor{gsvblue}{Durgesh Kumar Legha}}
\noindent
\textbf{Durgesh Kumar Legha} (Roll No.\ 2270007) is a final-year
B.Tech.\ student in Civil Engineering (Specialisation: Rail
Engineering) at Gati Shakti Vishwavidyalaya, Vadodara. His
undergraduate research focuses on railway-track substructure
behaviour, ballast degradation mechanisms, and sustainable
infrastructure for Indian Railways. He has co-authored two journal
manuscripts under review (\textit{Railway Engineering Science},
\textit{Construction and Building Materials}) and two accepted
Springer conference papers, and is a co-inventor on a pending Indian
utility-patent application for the large-scale drop-weight impact
apparatus described in this thesis.

\vspace{0.6cm}

\section*{\textcolor{gsvblue}{Sandeep Patel}}
\noindent
\textbf{Sandeep Patel} (Roll No.\ 2280043) is a final-year B.Tech.\
student in Civil Engineering (Specialisation: Rail Engineering) at
Gati Shakti Vishwavidyalaya, Vadodara. His undergraduate work
centres on the experimental geotechnics of railway ballast, including
the design, fabrication, and instrumentation of large-scale
laboratory apparatuses. He is a co-author on the
\textit{Construction and Building Materials} manuscript and on the
MedGU 2025 conference paper, and co-led the structured-light scanning
and per-particle morphological-analysis programme reported in
Chapters~\ref{chap:methodology} and~\ref{chap:morph-results} of this
thesis.

\vspace{1.0cm}
\section*{\textcolor{gsvblue}{Supervisor}}
\noindent
\textbf{Dr.\ Dinesh Gundavaram} is an Assistant Professor in the
School of Civil Engineering at Gati Shakti Vishwavidyalaya, Vadodara.
His research interests include railway geotechnology, cyclic-loading
behaviour of ballast and sub-ballast, resilient interlayer materials,
and the development of large-scale laboratory testing protocols
calibrated to Indian Railways field conditions. He is the
co-inventor on the pending Indian utility-patent application for the
apparatus reported in this thesis.

\vfill
\begin{center}
\rule{0.5\textwidth}{0.4pt}\\[6pt]
\textsc{\small Gati Shakti Vishwavidyalaya, Vadodara}\\
\textsc{\small April 2026}\\[6pt]
\rule{0.5\textwidth}{0.4pt}
\end{center}
